(1) Fix the orientation of $W_{1,3}$ given by the ordered basis $a_2, b_2, b_3, b_4$ of Section \ref{sec:setting}, written $e_0, \dots, e_3$ below. With boosts $B_i := e_0\wedge e_i$ and rotations $R_i := e_j\wedge e_k$ ($(i,j,k)$ cyclic), the $\mathfrak{so}(1,3)$ brackets are $[R_i, R_j] = \epsilon_{ijk}R_k$, $[R_i, B_j] = \epsilon_{ijk}B_k$, $[B_i, B_j] = -\epsilon_{ijk}R_k$ (the first two are rotations acting on vectors and axial pairs; the third is the boost-boost computation of \cite[Lem.~5.4]{ButlerT6}, reproved by the script). Setting $A^{(1)}_i := \tfrac12(R_i + iB_i)$, $A^{(2)}_i := \tfrac12(R_i - iB_i)$ and using $[B_i, R_j] = +\epsilon_{ijk}B_k$,
\begin{gather*}
[A^{(1)}_i, A^{(1)}_j] = \tfrac14\,\epsilon_{ijk}\bigl(R_k + iB_k + iB_k + R_k\bigr) = \epsilon_{ijk}A^{(1)}_k,\\
[A^{(1)}_i, A^{(2)}_j] = \tfrac14\,\epsilon_{ijk}\bigl(R_k - iB_k + iB_k - R_k\bigr) = 0,
\end{gather*}
and likewise $[A^{(2)}_i, A^{(2)}_j] = \epsilon_{ijk}A^{(2)}_k$; the script verifies all fifteen brackets of the $R, B$ table and the nine cross brackets exactly. Duality: $\ast B_i = -R_i$, $\ast R_i = B_i$ (Lemma \claimref{correspondence.lorentzian-duality-complex-structure}), so $\ast A^{(1)}_i = \tfrac12(B_i - iR_i) = -iA^{(1)}_i$ and $\ast A^{(2)}_i = +iA^{(2)}_i$: $\mathfrak{g}^+$ is spanned by the $A^{(2)}_i$ and $\mathfrak{g}^-$ by the $A^{(1)}_i$. Which eigenspace is called self-dual is a convention, and nothing below depends on it.
(2) (a) Two orderings of a word $q(\rho(P_1))\cdots q(\rho(P_r))$, $P_s \in \mathfrak{g}^-$, differ by a sequence of adjacent transpositions, each of which changes the word by the word with the adjacent pair $q(\rho(P_s))q(\rho(P_t))$ replaced by $[q(\rho(P_s)), q(\rho(P_t))] = \hbar\,q(\rho([P_s, P_t]))$ (Lemma \claimref{correspondence.symmetrised-quadratics-palatial-algebra}(2) and Proposition \claimref{correspondence.spin-representation}); since $[P_s, P_t] \in \mathfrak{g}^-$ by (1), the result is $\hbar$ times a word in elements of $q(\rho(\mathfrak{g}^-))$, an element of $\hbar\,\mathcal{A}^-$; for $r = 2$ it is the single commutator displayed. (b) For $P \in \mathfrak{g}^-$ and $Q \in \mathfrak{g}^+$, $[q(\rho(P)), q(\rho(Q))] = \hbar\,q(\rho([P,Q])) = 0$ by (1); the centraliser of $q(\rho(\mathfrak{g}^+))$ is a unital subalgebra containing the generators of $\mathcal{A}^-$, hence $\mathcal{A}^-$. (c) By Lemma \claimref{correspondence.symmetrised-quadratics-palatial-algebra}(1), $\operatorname{ad}_{q(\rho(Q))}/\hbar$ acts on the generators $Z^\alpha, \bar Z_\beta$ linearly, so it preserves the span $F_k$ of the words of length at most $k$; each $F_k$ is finite-dimensional, so $\mathcal{P} = \bigcup_k F_k$ is a locally finite $\mathfrak{g}^+$-module and Fact \claimref{correspondence.imported-complete-reducibility} decomposes each $F_k$, compatibly, into isotypic components; write $\mathcal{P} = \mathcal{P}^{\mathfrak{g}^+} \oplus \mathcal{P}_{\mathrm{ad}} \oplus \mathcal{P}_{\mathrm{rest}}$ for the trivial, adjoint and remaining isotypic parts. The map $Q \mapsto q(\rho(Q))$ intertwines the adjoint representation of $\mathfrak{g}^+$ with the adjoint action on $\mathcal{P}$ (Lemma \claimref{correspondence.symmetrised-quadratics-palatial-algebra}(2)) and is injective, so the googly slot is an irreducible subrepresentation isomorphic to the adjoint one, hence contained in $\mathcal{P}_{\mathrm{ad}}$. By (b), $\mathcal{A}^- \subseteq \mathcal{P}^{\mathfrak{g}^+}$, and $\mathcal{P}^{\mathfrak{g}^+} \cap (\mathcal{P}_{\mathrm{ad}} \oplus \mathcal{P}_{\mathrm{rest}}) = 0$; so every element of $\mathcal{A}^-$ has zero adjoint-isotypic component, whatever its degree.
